
\section{Aide-mémoire CS50}

\begin{UPSTIinfor}{style50 : Vérification de mise en page}
	Un code correct n'est pas simplement un code qui compile. \textbf{Il doit aussi et surtout être lisible et compréhensible par tous}
	\texttt{style50} est un outil vous permettant d'évaluer et d'améliorer la mise en page de votre code.

	Deux possibilités :
	\begin{itemize}
		\item Par l'interface graphique
		      \begin{itemize}
			      \item[$\Box$] En haut à droite de vsCode, cliquer sur \texttt{style50}.
		      \end{itemize}
		\item Par le terminal :
	\end{itemize}
\end{UPSTIinfor}
\begin{lstlisting}[language=bash,style=console]
style50 <nom_du_fichier>
\end{lstlisting}

\begin{UPSTIinfor}{Lancer une vérification automatique}
	\begin{itemize}
		\item[$\Box$] Vérifier que votre code passe  \texttt{style50} :
		      \begin{itemize}
			      \item[$\Box$] En haut à droite de vsCode, cliquer sur \texttt{style50}.
			      \item[$\Box$] Vous pouvez alors appliquer les changements suggérés en cliquant sur \texttt{Apply}
		      \end{itemize}
		\item[$\Box$] Vérifier que votre programme répond au cahier des charges :
		      \begin{lstlisting}[language=bash,style=console]
	check50 iut-geii-annecy/exercices/2026/info1/tp1/0_hello/variable
\end{lstlisting}
		      \begin{itemize}
			      \item[$\Box$] Si des erreurs sont détectées, corrigez-les et vérifiez à nouveau.
		      \end{itemize}
	\end{itemize}
\end{UPSTIinfor}